\documentclass[letterpaper,11pt]{article}

\usepackage[empty]{fullpage}
\usepackage[T1]{fontenc}
\usepackage{charter}
\usepackage[colorlinks=true, linkcolor=NavyBlue, urlcolor=NavyBlue]{hyperref}
\usepackage{parskip}

\addtolength{\oddsidemargin}{-0.5in}
\addtolength{\evensidemargin}{-0.5in}
\addtolength{\textwidth}{1.0in}
\addtolength{\topmargin}{-0.6in}
\addtolength{\textheight}{1.4in}

\setlength{\parskip}{6pt}

\begin{document}

\noindent\textbf{\large Charles Benello} \hfill \today \\
\noindent +1(312) 989-9950 $|$ cmbenello@uchicago.edu $|$ Chicago, IL

\vspace{8pt}

\noindent\textbf{Akuna Capital -- Options 201 Course Application}

\vspace{4pt}

I'm a Master's student in Financial Mathematics at UChicago, half a mile from your Loop office. I just finished the Options Pricing course this term. What pulled me toward this side of the field wasn't the math itself but the way the math has immediate, measurable consequences: get the Greeks wrong on a vol surface and you lose money in a way that getting a regression coefficient slightly off in a paper never punishes you for.

A snapshot of this year:

At PayPal, through UChicago's Project Lab, I'm a quantitative researcher working on transaction foundation models for fraud detection. I'm comparing BERT-style masked prediction against GPT-style next-token prediction across 24M+ transactions, with ablations on tokenization, time encoding, and context window size. The hard problem is calibrating probability on a 0.12\% positive rate, which is structurally similar to thinking about edge in low-probability tail events.

Last fall at Bank of America (also through Project Lab) I built an AI-driven platform to rebalance HQLA portfolios under capital, liquidity, and NII constraints, with a multi-agent LLM debate system that generates macro stress scenarios and produces probability-weighted reallocation recommendations. At Alter Domus, my current SWE internship, I lifted document extraction accuracy on 100+ page financial filings from a 50\% baseline to 92--97\% by architecting a multi-stage pipeline with cross-reference verification.

On the systems side: my paper CrocSort (VLDB 2026, co-first-author) builds a parallel external merge sort in Rust that completes within memory budgets where PostgreSQL, DuckDB, and ClickHouse abort. I sole-authored a follow-up at ADMS 2026 on sort configuration for cloud and disaggregated hardware. At EPFL last summer I cut out-of-memory query times by 50 to 200\% with a CUDA decompression strategy.

What I think sets me apart from a typical applicant is that I'm an engineer-first quant. Akuna's pitch is that competitive market-making is a software problem as much as it is a math problem, and I've spent three years on the software side of that equation, where a slow or wrong answer has measurable cost. I write performant Rust and C++ as comfortably as I write Python, I've built end-to-end ML pipelines, and I've shipped kernel-level CUDA. I also won the Wayne Booth Graduate Student Teaching Prize this year, UChicago's top graduate teaching award, for the ability to communicate hard concepts under time pressure to rooms of students each pulling for a different answer. I'd argue that's the closest thing in academic life to a trading desk.

The honest reason I want Options 201: I want to figure out whether trading is the right next step before committing. I've spent serious time inside research, full-stack engineering, and ML. I have not spent serious time inside a trading firm. A week of lectures, mini-shadows, and quizzes is the cleanest way for me to find out, and the fact that you're headquartered in Chicago makes the practical case obvious.

Thanks for considering me.

\vspace{4pt}

\noindent Charles Benello

\end{document}
